\section{Planteamiento del problema y justificación}

En tareas de transporte cooperativo, cada AGV dispone típicamente de percepción y comunicación limitadas, de modo que no resulta realista asumir conocimiento global perfecto del sistema. Al mismo tiempo, la carga compartida introduce acoplamientos dinámicos cuya influencia puede variar por cambios de masa aparente o por desplazamientos del centro de masas. Esta combinación dificulta mantener una trayectoria deseada y una formación compatible con la manipulación segura del objeto transportado.

La literatura clásica de consenso y control distribuido ofrece una base sólida para coordinar agentes con interacciones locales \citep{olfati2004consensus,ren2008distributed}; sin embargo, la presencia de incertidumbre inercial en la carga exige mecanismos adicionales de adaptación o robustez. La justificación académica del trabajo reside en comparar, en un marco controlado de simulación, el comportamiento de una ley nominal distribuida frente a una alternativa capaz de absorber incertidumbre paramétrica sin depender de hardware real ni de una infraestructura software compleja.
